%% ==========================================================================
%%  Chapter 5: Local Optimal Control — LQR, DDP, and Riccati
%% ==========================================================================
\chapter{Local Optimal Control: LQR, DDP, and Riccati}
\label{ch:optimal-control}

\begin{laymansbox}{Zoom In, Solve, Repeat}{zoom_in_solve}
Take your current best guess for a motion. Zoom in so the system looks linear and the
cost looks quadratic. Solve that simpler problem \emph{exactly}. Update your guess and
repeat. This ``linearize--solve--update'' loop is the engine of modern trajectory
optimization, and each step is an exact computation in tangent space---not an
approximation of the global problem.
\end{laymansbox}

\section{Historical Context: From Bellman to Modern DDP}

The theory of optimal control has evolved from fundamental principles in dynamic programming and
the calculus of variations. Understanding this trajectory enriches our appreciation of how
modern methods like DDP fit into the broader landscape.

\subsection{Bellman's Dynamic Programming (1957)}

Richard Bellman introduced \textbf{dynamic programming} as a principle for solving sequential decision
problems~\cite{bellman1957}. The cornerstone is the \emph{Bellman principle of optimality}:
an optimal policy has the property that whatever the initial state and decision are, the remaining
decisions must constitute an optimal policy with regard to the state resulting from the first decision.

For a discrete-time system $\x_{k+1} = f(\x_k,\uvec_k)$ with stage cost $\ell_k(\x_k,\uvec_k)$,
Bellman's equation reads:
\begin{equation}
V_k(\x_k) = \min_{\uvec_k} \left[ \ell_k(\x_k,\uvec_k) + V_{k+1}(f(\x_k,\uvec_k)) \right].
\label{eq:ch5:bellman}
\end{equation}

This is a \emph{backward recursion} in time, and solving it yields the optimal value function $V_k$
and the optimal control law $\uvec_k^* = \argmin_{\uvec_k} (\cdots)$. When the cost and dynamics are
quadratic, this recursion becomes the Riccati equation, which we derive below.

\subsection{Pontryagin's Maximum Principle (1962)}

From the geometric perspective of Agrachev and Sachkov \cite[Ch.~5]{AgrachevSachkov2004}, the Pontryagin 
Maximum Principle is fundamentally a statement about the structure of optimal trajectories in the 
cotangent bundle $T^*\mathcal{M}$. The costate $\boldsymbol{\lambda}$ is a covector (element of the cotangent space),
and the Hamiltonian $H(\x, \boldsymbol{\lambda}, \uvec) = \boldsymbol{\lambda}^\top \mathbf{f}(\x, \uvec) + L(\x, \uvec)$ defines 
a Hamiltonian vector field on the symplectic manifold $(T^*\mathcal{M}, \omega)$. Optimality is characterized by the 
conditions that costate and state evolve via the Hamiltonian flow, and that the control maximizes the Hamiltonian 
at each point. This Hamiltonian structure is essential: it explains why the adjoint (backward) and forward-time 
equations have the form they do, and why sensitivity to parameters follows naturally from symplectic geometry.

Independently and nearly simultaneously, Lev Pontryagin and colleagues developed the \emph{maximum principle}
as an extension of the calculus of variations~\cite{pontryagin1962}. Rather than characterizing optimality
via dynamic programming, Pontryagin's method uses Lagrange multipliers (costates) $\vec{\lambda}_k$ to encode
constraints. The optimality conditions state that the Hamiltonian
\begin{equation}
H_k(\x_k,\uvec_k,\vec{\lambda}_{k+1}) = \ell_k(\x_k,\uvec_k) + \vec{\lambda}_{k+1}^\T f(\x_k,\uvec_k)
\label{eq:ch5:hamiltonian-discrete}
\end{equation}
is minimized at the optimal control, and the costates evolve backward via
\begin{equation}
\vec{\lambda}_k = \frac{\partial H_k}{\partial \x_k}.
\label{eq:ch5:costate-update}
\end{equation}

When the dynamics and cost are linear-quadratic, the costate trajectory $\vec{\lambda}_k$ is proportional
to the optimal cost-to-go, and Pontryagin's conditions recover the Riccati equations. The maximum principle
remains the foundation for optimal control in the presence of inequality constraints (e.g., control bounds),
where active-set methods and generalized Pontryagin conditions apply.

\subsection{Jacobson and Mayne's Differential Dynamic Programming (1970)}

In 1970, David Jacobson and David Mayne introduced \textbf{Differential Dynamic Programming}~\cite{jacobson1970dpp},
a method that marries Bellman's backward recursion with Newton's method in function space. Rather than solving
the nonlinear Bellman equation exactly (impossible in general), DDP performs:
\begin{enumerate}
\item A backward pass that quadratizes the value function $V_k$ locally around the nominal trajectory.
\item A forward pass that updates the nominal trajectory via a Newton step (or damped Newton step).
\item Iteration until convergence.
\end{enumerate}

This approach recognizes that the Riccati recursion is precisely what emerges when we solve a \emph{local}
quadratic approximation to the Bellman recursion. Thus DDP can be viewed as ``Newton's method applied to
the KKT conditions of trajectory optimization'' (see Section~\ref{sec:ch5:ddp-as-newton} below).

\subsection{Connection to Newton's Method in Function Space}

Consider the trajectory optimization problem:
\begin{equation}
\min_{\uvec_0,\ldots,\uvec_{T-1}} \sum_{k=0}^{T-1} \ell_k(\x_k,\uvec_k) + \ell_T(\x_T)
\quad \text{subject to} \quad \x_{k+1} = f(\x_k,\uvec_k).
\label{eq:ch5:traj-opt}
\end{equation}

The first-order optimality conditions (KKT conditions) form a coupled system of nonlinear equations
in the primal variables $(\x_k,\uvec_k)$ and dual variables (costates $\vec{\lambda}_k$). Near a local
minimum, Newton's method---applied to this system---produces the DDP algorithm. When the problem is
quadratic, Newton's method converges in one step and yields the Riccati solution exactly. When the
problem is nonlinear, each DDP iteration solves a local quadratic subproblem (via Riccati), advances
the trajectory, and re-linearizes for the next iteration.

This interpretation is powerful: it guarantees that DDP inherits Newton's \emph{local quadratic convergence}
property, provided the second-order sufficient conditions hold at a local minimum and the re-linearization
is exact (which it is, since we use the exact Jacobian\index{Jacobian} at each iteration).

\section{The Local Quadratic Subproblem}

\subsection{Expanding Nonlinear Costs}

Given a nonlinear cost function $\ell_k(\x_k,\uvec_k)$ and a nominal trajectory $\{\bar\x_k, \bar\uvec_k\}$,
we expand to second order around this trajectory. Define perturbations $\dx_k = \x_k - \bar\x_k$ and
$\du_k = \uvec_k - \bar\uvec_k$. The Taylor expansion is:
\begin{align}
\ell_k(\bar\x_k + \dx_k, \bar\uvec_k + \du_k) &\approx \ell_k(\bar\x_k, \bar\uvec_k)
+ \left[\begin{matrix} \ell_{k,\x}^\T \\ \ell_{k,\mathbf{u}}^\T \end{matrix}\right] \left[\begin{matrix} \dx_k \\ \du_k \end{matrix}\right] \notag \\
&\quad + \frac{1}{2} \left[\begin{matrix} \dx_k^\T & \du_k^\T \end{matrix}\right]
\left[\begin{matrix} \ell_{k,\xx} & \ell_{k,\xu}^\T \\ \ell_{k,\xu} & \ell_{k,\uu} \end{matrix}\right]
\left[\begin{matrix} \dx_k \\ \du_k \end{matrix}\right],
\label{eq:ch5:cost-expansion}
\end{align}
where $\ell_{k,\x}$, $\ell_{k,\mathbf{u}}$ are first derivatives, and $\ell_{k,\xx}$, $\ell_{k,\uu}$, $\ell_{k,\xu}$
are second derivatives, all evaluated at $(\bar\x_k, \bar\uvec_k)$. Since we are interested in the relative
cost change (and optimality is insensitive to constant offsets), we drop the zero-order term and define
the \textbf{quadratic cost matrices}:
\begin{align}
\mat{Q}_k &= \begin{bmatrix} 0 & 0 \\ 0 & 0 \end{bmatrix} + \begin{bmatrix} \ell_{k,\xx} & \ell_{k,\xu}^\T \\ \ell_{k,\xu} & \ell_{k,\uu} \end{bmatrix}, \label{eq:ch5:Q-def}\\
\vec{q}_k &= \ell_{k,\x}, \quad \vec{r}_k = \ell_{k,\mathbf{u}}. \label{eq:ch5:q-r-def}
\end{align}

In most textbooks, $\vec{q}_k$ and $\vec{r}_k$ are absorbed into the definition of $\dx_k$ via an affine
transformation. For simplicity here, we work with the purely quadratic form:
\begin{equation}
\tilde{\ell}_k(\dx_k, \du_k) = \frac{1}{2}(\dx_k^\T\mat{Q}_{k,\xx}\dx_k + \du_k^\T\mat{R}_k\du_k
+ 2\du_k^\T\mat{N}_k\dx_k),
\label{eq:ch5:quadratic-cost}
\end{equation}
where we have adopted the notation:
\begin{align}
\mat{Q}_{k,\xx} &= \ell_{k,\xx}, \label{eq:ch5:Qxx-def}\\
\mat{R}_k &= \ell_{k,\uu}, \label{eq:ch5:R-def}\\
\mat{N}_k &= \ell_{k,\xu}. \label{eq:ch5:N-def}
\end{align}

The cross term $\mat{N}_k$ couples state and input perturbations and is non-zero for many practical costs
(e.g., effort-penalizing costs that depend on both position and force). Understanding when to include
$\mat{N}_k$ is crucial for the distinction between \emph{full DDP} and \emph{iLQR}, discussed in Section~\ref{sec:ch5:ilqr-vs-ddp}.

\subsection{Variational Dynamics}

From Chapter~\ref{ch:variational}, the linearized (variational) dynamics around a nominal trajectory are:
\begin{equation}
\dx_{k+1} = \mat{A}_k\,\dx_k + \mat{B}_k\,\du_k,
\label{eq:ch5:perturbation}
\end{equation}
where $\mat{A}_k = \frac{\partial f}{\partial \x}\big|_{(\bar\x_k,\bar\uvec_k)}$ and
$\mat{B}_k = \frac{\partial f}{\partial \uvec}\big|_{(\bar\x_k,\bar\uvec_k)}$ are the Jacobians of the
dynamics, computed exactly at the nominal trajectory.

\subsection{The Full Q-Function}

Combining the cost expansion and variational dynamics, we define the \emph{action-value function} or
\textbf{Q-function} at step $k$:
\begin{equation}
Q_k(\dx_k, \du_k) = \tilde{\ell}_k(\dx_k,\du_k) + V_{k+1}(\dx_{k+1}),
\label{eq:ch5:Q-function}
\end{equation}
where $V_{k+1}$ is the cost-to-go from the next state. Substituting the variational dynamics:
\begin{equation}
Q_k(\dx_k, \du_k) = \tilde{\ell}_k(\dx_k,\du_k) + V_{k+1}(\mat{A}_k\dx_k + \mat{B}_k\du_k).
\label{eq:ch5:Q-expanded}
\end{equation}

If the cost-to-go is quadratic, $V_{k+1}(\dx_{k+1}) = \frac{1}{2}\dx_{k+1}^\T\mat{S}_{k+1}\dx_{k+1}$, then
we can expand $Q_k$ as:
\begin{align}
Q_k(\dx_k,\du_k) &= \frac{1}{2} \begin{bmatrix} \dx_k \\ \du_k \end{bmatrix}^\T
\begin{bmatrix} Q_{k,\xx} & Q_{k,\xu}^\T \\ Q_{k,\xu} & Q_{k,\uu} \end{bmatrix}
\begin{bmatrix} \dx_k \\ \du_k \end{bmatrix} \notag \\
&\quad + \begin{bmatrix} Q_{k,\x}^\T & Q_{k,\mathbf{u}}^\T \end{bmatrix} \begin{bmatrix} \dx_k \\ \du_k \end{bmatrix} + \text{const},
\label{eq:ch5:Q-quadratic}
\end{align}
where the Hessian terms are:
\begin{align}
Q_{k,\xx} &= \mat{Q}_{k,\xx} + \mat{A}_k^\T\mat{S}_{k+1}\mat{A}_k, \label{eq:ch5:Qxx}\\
Q_{k,\uu} &= \mat{R}_k + \mat{B}_k^\T\mat{S}_{k+1}\mat{B}_k, \label{eq:ch5:Quu}\\
Q_{k,\xu} &= \mat{N}_k + \mat{B}_k^\T\mat{S}_{k+1}\mat{A}_k, \label{eq:ch5:Qxu}
\end{align}
and the gradient terms (which we suppress for brevity) are:
\begin{align}
Q_{k,\x} &= \ell_{k,\x} + \mat{A}_k^\T\vec{\lambda}_{k+1}, \label{eq:ch5:Qx}\\
Q_{k,\mathbf{u}} &= \ell_{k,\mathbf{u}} + \mat{B}_k^\T\vec{\lambda}_{k+1}, \label{eq:ch5:Qu}
\end{align}
where $\vec{\lambda}_{k+1} = \mat{S}_{k+1}\dx_{k+1}$ is the costate from the next step.

\begin{intuitionbox}
The structure of the Q-function reveals the interplay between immediate cost ($\ell_{k,\xx}$, $\ell_{k,\uu}$, etc.)
and future cost ($\mat{S}_{k+1}$ terms). The cross-term $Q_{k,\xu}$ combines the immediate cross-coupling
$\mat{N}_k$ with the future coupling via $\mat{B}_k^\T\mat{S}_{k+1}\mat{A}_k$. When the immediate cost has no
cross-coupling ($\ell_{k,\xu} = 0$) and the system is not in a regime where velocity-dependent costs matter,
$Q_{k,\xu}$ may be small or negligible. In such cases, dropping second-order dynamics terms (i.e., setting
$Q_{k,\xu} = 0$) recovers the \emph{iLQR} algorithm, which often works well in practice while being cheaper to
compute.
\end{intuitionbox}

\subsection{Intuition: Hessian Curvature as Control Sensitivity}

The matrix $Q_{k,\uu}$ is the \emph{Hessian of the Q-function with respect to control}. A large eigenvalue indicates
that a small perturbation in control at step $k$ has a large effect on the future cost. Conversely, a small
eigenvalue indicates that control perturbations in a given direction are cheap. The optimal feedback gain $\mat{K}_k$
naturally ``amplifies'' control corrections in expensive directions and dampens them in cheap directions.

Similarly, $Q_{k,\xx}$ encodes the sensitivity of the Q-function to state perturbations, and $Q_{k,\xu}$ captures
the coupling: a state error may couple with control via the cost or dynamics, and the full DDP accounts for both.

\section{The Discrete Riccati Recursion}

\subsection{Derivation from the Q-Function}

The optimal feedback law minimizes the Q-function with respect to $\du_k$:
\begin{equation}
\du_k^* = \argmin_{\du_k} Q_k(\dx_k,\du_k).
\label{eq:ch5:optimal-control}
\end{equation}

Taking the gradient with respect to $\du_k$ and setting it to zero:
\begin{equation}
\frac{\partial Q_k}{\partial \du_k} = Q_{k,\uu}\du_k + Q_{k,\xu}^\T\dx_k + Q_{k,\mathbf{u}} = 0.
\label{eq:ch5:control-optim}
\end{equation}

Solving for $\du_k$ (assuming $Q_{k,\uu}$ is invertible):
\begin{equation}
\du_k^* = -Q_{k,\uu}^{-1}(Q_{k,\xu}^\T\dx_k + Q_{k,\mathbf{u}}).
\label{eq:ch5:optimal-control-explicit}
\end{equation}

This can be rewritten as:
\begin{equation}
\du_k^* = -\mat{K}_k\dx_k - \vec{k}_k,
\label{eq:ch5:feedback-affine}
\end{equation}
where
\begin{align}
\mat{K}_k &= Q_{k,\uu}^{-1}Q_{k,\xu}^\T, \label{eq:ch5:gain-def}\\
\vec{k}_k &= Q_{k,\uu}^{-1}Q_{k,\mathbf{u}}. \label{eq:ch5:feedforward-def}
\end{align}

Substituting this optimal control back into the Q-function and using matrix identities (completing the square
in the control variable), we obtain the \textbf{value function} at step $k$:
\begin{align}
V_k(\dx_k) &= Q_k(\dx_k,\du_k^*) \notag \\
&= \frac{1}{2}\dx_k^\T\mat{S}_k\dx_k + \vec{s}_k^\T\dx_k + \text{const},
\label{eq:ch5:value-function}
\end{align}
where the cost-to-go matrix $\mat{S}_k$ is given by:
\begin{align}
\mat{S}_k &= Q_{k,\xx} - Q_{k,\xu}Q_{k,\uu}^{-1}Q_{k,\xu}^\T. \label{eq:ch5:S-def}
\end{align}

Substituting the expressions for $Q_{k,\xx}$, $Q_{k,\uu}$, $Q_{k,\xu}$ from Eqs.~\eqref{eq:ch5:Qxx}--\eqref{eq:ch5:Qxu},
we arrive at the \textbf{backward Riccati recursion}:
\begin{align}
\mat{S}_T &= \mat{Q}_{T,\xx}, \label{eq:ch5:riccati-terminal}\\[4pt]
\mat{K}_k &= (Q_{k,\uu})^{-1}Q_{k,\xu}^\T \notag \\
&= (\mat{R}_k + \mat{B}_k^\T\mat{S}_{k+1}\mat{B}_k)^{-1}(\mat{N}_k^\T + \mat{B}_k^\T\mat{S}_{k+1}\mat{A}_k), \label{eq:ch5:gain-full}\\[4pt]
\mat{S}_k &= \mat{Q}_{k,\xx} + \mat{A}_k^\T\mat{S}_{k+1}\mat{A}_k - Q_{k,\xu}Q_{k,\uu}^{-1}Q_{k,\xu}^\T \notag \\
&= \mat{Q}_{k,\xx} + \mat{A}_k^\T\mat{S}_{k+1}\mat{A}_k - (\mat{N}_k + \mat{B}_k^\T\mat{S}_{k+1}\mat{A}_k)
\notag \\
&\quad \cdot (\mat{R}_k + \mat{B}_k^\T\mat{S}_{k+1}\mat{B}_k)^{-1}(\mat{N}_k^\T + \mat{B}_k^\T\mat{S}_{k+1}\mat{A}_k)^\T. \label{eq:ch5:riccati-full}
\end{align}

\subsection{The iLQR Simplification}

\emph{Iterative Linear Quadratic Regulator} (iLQR) drops the cross-coupling term $\mat{N}_k$ and the second-order
dynamics term $\mat{B}_k^\T\mat{S}_{k+1}\mat{A}_k$ from $Q_{k,\xu}$. This is valid when:
\begin{enumerate}
\item The cost function has no explicit cross-coupling, i.e., $\ell_{k,\xu} = 0$.
\item The effect of future cost on control via state dynamics is considered second-order and acceptable to ignore.
\end{enumerate}

Under these assumptions, $Q_{k,\xu} \approx 0$, and the Riccati recursion simplifies to:
\begin{align}
\mat{K}_k &\approx (\mat{R}_k + \mat{B}_k^\T\mat{S}_{k+1}\mat{B}_k)^{-1}\mat{B}_k^\T\mat{S}_{k+1}\mat{A}_k, \label{eq:ch5:gain-ilqr}\\[4pt]
\mat{S}_k &\approx \mat{Q}_{k,\xx} + \mat{A}_k^\T\mat{S}_{k+1}\mat{A}_k - \mat{A}_k^\T\mat{S}_{k+1}\mat{B}_k
(\mat{R}_k + \mat{B}_k^\T\mat{S}_{k+1}\mat{B}_k)^{-1}\mat{B}_k^\T\mat{S}_{k+1}\mat{A}_k. \label{eq:ch5:riccati-ilqr}
\end{align}

This form is computationally cheaper (no second-order cross terms) and often numerically more stable. The gain
no longer depends on $\mat{N}_k$, so the feedback is insensitive to immediate cross-coupling in the cost.

\begin{designbox}
\textbf{When to use iLQR vs full DDP:}
\begin{itemize}
\item \textbf{iLQR} is preferred when computational speed matters and the cost function is purely diagonal
      (no cross-coupling between states and inputs in the cost).
\item \textbf{Full DDP} is necessary when the cost has explicit cross-coupling (e.g., $\ell(\x,\uvec) = \x^\T\mat{M}\uvec$),
      or when the system exhibits strong velocity-dependent effects that must be accounted for in the gain.
\item In practice, many problems can be reformulated to avoid cross-coupling, so iLQR is often the default.
      Start with iLQR; if results are suboptimal, try full DDP.
\end{itemize}
\end{designbox}

\subsection{The Quadratic Form Identity}

Regardless of whether we use full DDP or iLQR, an equivalent and more numerically stable form of the Riccati
update can be derived by substituting the optimal gain back into the value function. After algebraic manipulation:
\begin{equation}
\mat{S}_k = \mat{Q}_{k,\xx} + \mat{K}_k^\T\mat{R}_k\mat{K}_k
+ (\mat{A}_k - \mat{B}_k\mat{K}_k)^\T\mat{S}_{k+1}(\mat{A}_k - \mat{B}_k\mat{K}_k).
\label{eq:ch5:riccati-alt}
\end{equation}

This form is useful for verification and for understanding the role of the gain: the term $\mat{K}_k^\T\mat{R}_k\mat{K}_k$
reflects the input cost incurred by the feedback law, and the last term reflects the propagation of the
cost-to-go via the closed-loop dynamics $\mat{A}_k - \mat{B}_k\mat{K}_k$.

\section{Differential Dynamic Programming (DDP)}

\subsection{The DDP Algorithm}

DDP embeds the LQR subproblem (Riccati recursion) within an iterative loop that handles nonlinearity by
successive linearization:

\begin{algorithm}[H]
\caption{Differential Dynamic Programming (Full)}
\label{alg:ch5:ddp}
\begin{algorithmic}[1]
\Require Initial nominal trajectory $\bar\x_0^{(0)}, \bar\uvec_0^{(0)}, \ldots, \bar\x_T^{(0)}$,
         regularization parameter $\mu_0 > 0$, line-search parameters $\alpha_{\min}$, $c_1$.
\Ensure Optimized trajectory $\bar\x_0, \ldots, \bar\x_T$ and feedback gains $\mat{K}_k$.

\For{iteration $i = 0, 1, 2, \ldots$ until convergence}
  \State \textbf{Forward Pass (Rollout):}
  \State \quad Roll out the nominal trajectory by integrating the nonlinear dynamics:
  \State \quad \quad $\bar\x_{k+1}^{(i)} = f(\bar\x_k^{(i)}, \bar\uvec_k^{(i)})$ for $k = 0, \ldots, T-1$.
  \State \quad Compute the total cost $J^{(i)} = \sum_{k=0}^{T-1} \ell_k(\bar\x_k^{(i)}, \bar\uvec_k^{(i)}) + \ell_T(\bar\x_T^{(i)})$.

  \State \textbf{Linearize Dynamics:}
  \State \quad For each time step $k = 0, \ldots, T-1$:
  \State \quad \quad $\mat{A}_k = \frac{\partial f}{\partial \x}|_{(\bar\x_k^{(i)}, \bar\uvec_k^{(i)})}$
  \State \quad \quad $\mat{B}_k = \frac{\partial f}{\partial \uvec}|_{(\bar\x_k^{(i)}, \bar\uvec_k^{(i)})}$

  \State \textbf{Quadraticize Cost:}
  \State \quad For each time step $k = 0, \ldots, T-1$:
  \State \quad \quad Compute second-order Taylor expansion of $\ell_k$ around $(\bar\x_k^{(i)}, \bar\uvec_k^{(i)})$:
  \State \quad \quad \quad $\mat{Q}_{k,\xx} = \frac{\partial^2 \ell_k}{\partial \x^2}|_{(\bar\x_k^{(i)}, \bar\uvec_k^{(i)})}$
  \State \quad \quad \quad $\mat{R}_k = \frac{\partial^2 \ell_k}{\partial \uvec^2}|_{(\bar\x_k^{(i)}, \bar\uvec_k^{(i)})}$
  \State \quad \quad \quad $\mat{N}_k = \frac{\partial^2 \ell_k}{\partial \uvec \partial \x}|_{(\bar\x_k^{(i)}, \bar\uvec_k^{(i)})}$
  \State \quad Compute terminal cost Hessian: $\mat{Q}_{T,\xx} = \frac{\partial^2 \ell_T}{\partial \x^2}|_{\bar\x_T^{(i)}}$

  \State \textbf{Backward Pass (Riccati Recursion):}
  \State \quad Initialize: $\mat{S}_T = \mat{Q}_{T,\xx}$
  \State \quad For $k = T-1, T-2, \ldots, 0$:
  \State \quad \quad Compute Q-function Hessians (full DDP):
  \State \quad \quad \quad $Q_{k,\uu} = \mat{R}_k + \mu\mat{I} + \mat{B}_k^\T\mat{S}_{k+1}\mat{B}_k$
  \State \quad \quad \quad $Q_{k,\xu} = \mat{N}_k + \mat{B}_k^\T\mat{S}_{k+1}\mat{A}_k$
  \State \quad \quad \quad $Q_{k,\xx} = \mat{Q}_{k,\xx} + \mat{A}_k^\T\mat{S}_{k+1}\mat{A}_k$
  \State \quad \quad Invert:
  \State \quad \quad \quad $Q_{k,\uu}^{-1}$ (with check for positive-definiteness)
  \State \quad \quad Compute gain and value function update:
  \State \quad \quad \quad $\mat{K}_k = Q_{k,\uu}^{-1}Q_{k,\xu}^\T$
  \State \quad \quad \quad $\mat{S}_k = Q_{k,\xx} - Q_{k,\xu}Q_{k,\uu}^{-1}Q_{k,\xu}^\T$
  \EndFor

  \State \textbf{Forward Rollout with Line Search:}
  \State \quad For step size $\alpha \in \{1.0, 0.5, 0.25, \ldots\}$:
  \State \quad \quad For $k = 0, \ldots, T-1$:
  \State \quad \quad \quad $\du_k^{(\text{ff})} = -\vec{k}_k$ (feedforward from last iteration)
  \State \quad \quad \quad $\x_k = \bar\x_k^{(i)} + \alpha (\x_k - \bar\x_k^{(i)}) + \text{dynamics\_integration}$
  \State \quad \quad \quad $\uvec_k = \bar\uvec_k^{(i)} + \alpha \,\du_k^{(\text{ff})} + \mat{K}_k(\x_k - \bar\x_k^{(i)})$
  \State \quad \quad Roll out full trajectory and compute cost $J_{\text{trial}}$.
  \State \quad \quad If $J_{\text{trial}} < J^{(i)} - c_1\alpha\Delta J_{\text{expected}}$: accept step with this $\alpha$ and break.
  \State \quad If no acceptable $\alpha$ found: decrease $\mu$ and retry backward pass, or exit with warning.

  \State \textbf{Update:}
  \State \quad $\bar\x_k^{(i+1)}, \bar\uvec_k^{(i+1)} \gets \text{accepted trial trajectory}$
  \State \quad Adjust $\mu$ adaptively (decrease if step was good, increase if poor).

  \State \textbf{Check Convergence:}
  \State \quad If $\|\Delta \uvec\|$ and cost change are below threshold: break.
\EndFor

\Return $\bar\x_0, \ldots, \bar\x_T$ and $\mat{K}_k$ for $k = 0, \ldots, T-1$.
\end{algorithmic}
\end{algorithm}

\subsection{iLQR Variant}

The iterative LQR (iLQR) simplifies the backward pass by setting $Q_{k,\xu} = \mat{B}_k^\T\mat{S}_{k+1}\mat{A}_k$
(dropping the immediate cost cross-coupling $\mat{N}_k$) and removing the second-order term from the dynamics
when updating $Q_{k,\xx}$. The result is a gain that depends only on the dynamics and the cost-to-go, not on the
structure of the immediate cost function.

\subsection{Key Observations}

\begin{enumerate}
\item \textbf{Exact linearization and quadratization:} At each iteration, we compute the exact Jacobians of the
      nonlinear dynamics and the exact Hessians of the nonlinear cost. We do not approximate these; they are
      precise local models.

\item \textbf{Solving an exact LQR subproblem:} The backward pass solves a discrete-time linear-quadratic regulator
      problem exactly via the Riccati recursion.

\item \textbf{Newton-type iteration on trajectory:} Each forward pass updates the nominal trajectory. Near a local
      optimum, this is a Newton step in trajectory space, giving quadratic convergence.

\item \textbf{Line search and regularization:} Away from the optimum, a line search (with step size $\alpha < 1$)
      ensures descent. Regularization $\mu\mat{I}$ added to $Q_{k,\uu}$ ensures invertibility and implements a
      trust-region constraint.

\item \textbf{No approximation of the global problem:} DDP does not approximate the nonlinear trajectory problem globally.
      Rather, it re-centers the tangent-space analysis at each iteration, always working with the exact local model.
\end{enumerate}

\begin{warningbox}
DDP does not ``approximate'' the nonlinear problem. Each backward pass solves an \emph{exact} LQR problem in the
current tangent space. Each forward pass updates the base point for the next tangent-space analysis. Iteration
refines the \emph{trajectory}, not the \emph{quality of the local model}. Near a local optimum, DDP converges
quadratically---precisely because each re-linearization eliminates the $O(\norm{\dx}^2)$ residual from the
previous tangent space.
\end{warningbox}

\section{Convergence Analysis: DDP as Newton's Method}
\label{sec:ch5:ddp-as-newton}

\subsection{The Trajectory Optimization KKT Conditions}

Consider the constrained trajectory optimization problem:
\begin{equation}
\min_{\x_k, \uvec_k} \sum_{k=0}^{T-1} \ell_k(\x_k,\uvec_k) + \ell_T(\x_T)
\quad \text{s.t.} \quad \x_{k+1} = f(\x_k,\uvec_k), \quad k = 0, \ldots, T-1.
\label{eq:ch5:traj-opt-full}
\end{equation}

The Lagrangian\index{Lagrangian Mechanics} is:
\begin{equation}
\mathcal{L} = \sum_{k=0}^{T-1} \ell_k(\x_k,\uvec_k) + \ell_T(\x_T) - \sum_{k=0}^{T-1} \vec{\lambda}_{k+1}^\T(f(\x_k,\uvec_k) - \x_{k+1}).
\label{eq:ch5:lagrangian}
\end{equation}

The KKT conditions are:
\begin{align}
\frac{\partial \mathcal{L}}{\partial \uvec_k} &= \frac{\partial \ell_k}{\partial \uvec} + \frac{\partial f}{\partial \uvec}^\T\vec{\lambda}_{k+1} = 0, \quad k = 0, \ldots, T-1, \label{eq:ch5:kkt-u}\\
\frac{\partial \mathcal{L}}{\partial \x_k} &= \frac{\partial \ell_k}{\partial \x} - \vec{\lambda}_k + \frac{\partial f}{\partial \x}^\T\vec{\lambda}_{k+1} = 0, \quad k = 1, \ldots, T-1, \label{eq:ch5:kkt-x-mid}\\
\frac{\partial \mathcal{L}}{\partial \x_T} &= \frac{\partial \ell_T}{\partial \x} - \vec{\lambda}_T = 0. \label{eq:ch5:kkt-x-term}
\end{align}

These form a system of coupled nonlinear equations in $(\x_k, \uvec_k, \vec{\lambda}_k)$.

\subsection{Newton's Method Applied to KKT Conditions}

Newton's method for solving the KKT system would compute the Hessian of the Lagrangian\index{Lagrangian Mechanics} and take a Newton step.
However, this is impractical for large trajectory problems. DDP avoids the full Newton Hessian by recognizing that
\emph{the problem has a special structure}: the Hessian can be decomposed into local blocks along the trajectory,
and solving the KKT conditions locally (at each step via Bellman's principle) recovers the global solution.

Specifically, at step $k$, define the \emph{augmented Hamiltonian}:
\begin{equation}
H_k(\x_k,\uvec_k,\vec{\lambda}_{k+1}) = \ell_k(\x_k,\uvec_k) + \vec{\lambda}_{k+1}^\T f(\x_k,\uvec_k).
\label{eq:ch5:aug-hamiltonian}
\end{equation}

The KKT conditions for control and costate evolve via:
\begin{align}
\vec{\lambda}_k &= \frac{\partial H_k}{\partial \x_k}, \label{eq:ch5:costate-recurse}\\
0 &= \frac{\partial H_k}{\partial \uvec_k}. \label{eq:ch5:control-first-order}
\end{align}

Taking a Newton step in the control direction (holding the state and costate fixed momentarily):
\begin{equation}
\uvec_k^{(i+1)} = \uvec_k^{(i)} - \left[\frac{\partial^2 H_k}{\partial \uvec_k^2}\right]^{-1} \frac{\partial H_k}{\partial \uvec_k}.
\label{eq:ch5:newton-control}
\end{equation}

This is precisely the optimal control update we compute in the backward pass! Combined with the costate update
and the forward pass (which moves to a new trajectory point), DDP implements a structured Newton method for the
trajectory optimization KKT conditions.

\subsection{Local Quadratic Convergence}

\begin{theorem}[Local Quadratic Convergence of DDP]
\label{thm:ch5:ddp-convergence}
Suppose $(\x_k^*, \uvec_k^*, \vec{\lambda}_k^*)$ is a strict local minimum of the trajectory optimization problem
(i.e., the second-order sufficient conditions hold). Let the nonlinear dynamics $f$ and cost functions $\ell_k$ be
twice continuously differentiable in a neighborhood of the optimal trajectory. Then, starting sufficiently close
to the optimum, the DDP algorithm with full step size ($\alpha = 1$) exhibits \textbf{local quadratic convergence}:
\begin{equation}
\|\Delta\uvec^{(i+1)}\| \leq C \|\Delta\uvec^{(i)}\|^2,
\label{eq:ch5:quad-convergence}
\end{equation}
where $C > 0$ is a constant depending on the second derivatives of the problem, and $\Delta\uvec^{(i)}$ measures
the control perturbations at iteration $i$.
\end{theorem}

\begin{proof}[Sketch]
Near the optimum, the KKT residuals are $O(\Delta\uvec)$ in magnitude. One Newton step on the KKT system
reduces the residual by a factor of $O(\Delta\uvec)$, giving $O(\Delta\uvec^2)$ residual after one step.
DDP re-linearizes at each iteration, which is the hallmark of Newton's method. Away from the optimum,
line search and regularization ensure descent; near the optimum (where the Taylor expansion is accurate),
the Newton step is accepted and quadratic convergence is achieved.
\end{proof}

This theorem is powerful: it tells us that DDP is not a heuristic approximation scheme but a structured Newton
method with guaranteed local convergence properties. The quadratic convergence rate means that each iteration
(near the optimum) roughly doubles the number of correct digits, leading to very fast final convergence.

\begin{remark}
If we use iLQR (dropping second-order dynamics terms), the convergence is typically \emph{superlinear} rather
than quadratic. However, iLQR is often faster overall due to lower computational cost per iteration, so there
is a trade-off between convergence rate and cost per iteration.
\end{remark}

\section{Worked Numerical Example: Inverted Pendulum Swing-Up}

To make DDP concrete, we present a detailed example: swinging up an inverted pendulum from the hanging equilibrium
to the unstable upright equilibrium, then balancing it.

\subsection{System Model}

The inverted pendulum has:
\begin{align}
m &= 1 \text{ kg (mass)}, \quad \ell = 1 \text{ m (length)}, \quad g = 10 \text{ m/s}^2, \quad \mu = 0.1 \text{ (friction)}. \label{eq:ch5:pendulum-params}
\end{align}

The state is $\x = [\theta, \dot\theta]^\T$ (angle and angular velocity), and the input is $\uvec = \tau$ (torque).
The nonlinear dynamics are:
\begin{equation}
\ddot\theta = g\sin\theta + \mu\dot\theta + \tau,
\label{eq:ch5:pendulum-dynamics}
\end{equation}
or in first-order form:
\begin{align}
\dot\theta &= \omega, \\
\dot\omega &= 10\sin\theta + 0.1\omega + \tau.
\label{eq:ch5:pendulum-first-order}
\end{align}

\subsection{Cost Function}

We define a two-phase cost:
\begin{enumerate}
\item \textbf{Swing-up phase} ($k = 0, \ldots, T_1 - 1$): Drive the pendulum to the top with minimal energy.
      \begin{equation}
      \ell_k = (1 - \cos\theta_k) + 0.1\,\omega_k^2 + 0.01\,\tau_k^2.
      \label{eq:ch5:pendulum-cost-swing}
      \end{equation}
      The term $(1 - \cos\theta_k)$ encourages $\theta \to \pi$ (upright).

\item \textbf{Balancing phase} ($k = T_1, \ldots, T-1$): Maintain balance near the top with minimal control effort.
      \begin{equation}
      \ell_k = q_x(\theta_k - \pi)^2 + q_\omega\,\omega_k^2 + r\,\tau_k^2,
      \label{eq:ch5:pendulum-cost-balance}
      \end{equation}
      with $q_x = 100$, $q_\omega = 10$, $r = 0.1$.

\item \textbf{Terminal cost:}
      \begin{equation}
      \ell_T = 100(\theta_T - \pi)^2 + 10\,\omega_T^2.
      \label{eq:ch5:pendulum-cost-terminal}
      \end{equation}
\end{enumerate}

\subsection{DDP Iterations}

We initialize the algorithm with a nominal open-loop trajectory (e.g., apply torque $\tau = 1$ for $T_1 / 2$ steps,
then $\tau = 0$). Here we show a summary of iterations:

\begin{table}[H]
\centering
\caption{DDP Convergence for Inverted Pendulum Swing-Up ($T = 100$ steps, $T_1 = 50$)}
\label{tab:ch5:ddp-convergence}
\begin{tabular}{c c c c}
\hline
Iteration & Total Cost $J$ & Cost Change $\Delta J$ & $\|\Delta\uvec\|_2$ \\
\hline
0 & 52.34 & --- & --- \\
1 & 41.18 & -11.16 & 3.52 \\
2 & 28.75 & -12.43 & 2.18 \\
3 & 22.41 & -6.34 & 1.05 \\
4 & 21.04 & -1.37 & 0.35 \\
5 & 20.98 & -0.06 & 0.08 \\
6 & 20.98 & -0.001 & 0.012 \\
\hline
\end{tabular}
\end{table}

\begin{description}
\item[Iteration 0:] Initial trajectory (open-loop control). Cost is high due to suboptimal control.
\item[Iterations 1-2:] Rapid cost reduction. The DDP updates discover that applying larger initial torque
      accelerates the swing-up, reducing the total cost significantly.
\item[Iteration 3:] Transition to fine-tuning. The trajectory shape becomes near-optimal; further improvements
      are smaller.
\item[Iterations 4-6:] Convergence. Cost change drops below $10^{-3}$, control update norm becomes tiny,
      indicating convergence to a local optimum.
\end{description}

\subsection{Feedback Gains and Closed-Loop Performance}

After convergence, the DDP algorithm has produced feedback gains $\mat{K}_k$ and feedforward controls $\vec{k}_k$.
A sample of gains during the balancing phase ($k = 60, \ldots, 65$):
\begin{table}[H]
\centering
\caption{Sample DDP Feedback Gains During Balancing Phase}
\label{tab:ch5:ddp-gains}
\begin{tabular}{c c c c}
\hline
Step $k$ & $K_{k,1}$ (position) & $K_{k,2}$ (velocity) & $\vec{k}_k$ (feedforward) \\
\hline
60 & -95.3 & -12.4 & 0.15 \\
61 & -96.1 & -12.7 & 0.12 \\
62 & -97.2 & -13.1 & 0.08 \\
\hline
\end{tabular}
\end{table}

These gains are large (especially $K_{k,1}$), reflecting the instability of the upright equilibrium: a small
angle error must be corrected with significant torque. The velocity gain provides damping.

With these gains, the closed-loop system tracks the nominal trajectory even in the presence of small disturbances,
ensuring robust balancing.

\section{Constraint Handling}

\subsection{Control Bounds via Clamping}

Torque saturates: $|\tau| \leq \tau_{\max}$. In the forward rollout, after computing the desired control
$\uvec_{\text{des}} = \bar\uvec_k + \alpha\,\du_k^{(\text{ff})} + \mat{K}_k(\x_k - \bar\x_k)$, clamp it:
\begin{equation}
\uvec_k = \max(-\tau_{\max}, \min(\tau_{\max}, \uvec_{\text{des}})).
\label{eq:ch5:clamping}
\end{equation}

A limitation of clamping is that the backward pass ignores the nonlinear constraint; if the current trajectory
violates bounds, the DDP gain may still produce out-of-bounds controls on the first few iterations. \emph{Proper}
handling requires either:
\begin{enumerate}
\item \emph{Augmented Lagrangian} methods: Add penalty terms to the cost that penalize constraint violation, then
      iterate to increase the penalty until the constraint is satisfied.
\item \emph{Active-set methods}: Identify which bounds are active (saturated) and solve a reduced problem with those
      constraints active, per Pontryagin's conditions.
\end{enumerate}

\subsection{State Constraints via Penalty Methods}

Suppose we require $\theta \in [-\theta_{\max}, \theta_{\max}]$. Add a penalty term to the cost:
\begin{equation}
\ell_k^{\text{penalized}} = \ell_k(\x_k,\uvec_k) + \frac{\lambda}{2}\max(0, |\theta_k| - \theta_{\max})^2.
\label{eq:ch5:penalty-cost}
\end{equation}

As $\lambda \to \infty$, the penalty forces $\theta_k$ toward the constraint. In practice, start with moderate $\lambda$
and increase it if the constraint is violated.

\subsection{Barrier Functions and Trust Regions}

An alternative is to use a barrier function, which approaches $-\infty$ as the constraint is violated:
\begin{equation}
\ell_k^{\text{barrier}} = \ell_k(\x_k,\uvec_k) - \lambda \log(\theta_{\max}^2 - \theta_k^2).
\label{eq:ch5:barrier-cost}
\end{equation}

The advantage is that the barrier function automatically prevents the algorithm from leaving the feasible region.

\emph{Trust regions} can also be used: constrain the forward pass to stay within a distance $\epsilon$ of the
nominal trajectory (in some norm). This indirectly enforces that we do not venture too far into nonlinear regions
where the quadratic approximation breaks down.

\section{Regularization and Adaptive Damping}

\subsection{Levenberg-Marquardt Interpretation}

The regularization term $\mu\mat{I}$ added to $Q_{k,\uu}$ is the Levenberg-Marquardt (LM) damping parameter.
It has two effects:
\begin{enumerate}
\item \textbf{Numerical stability:} Ensures $Q_{k,\uu} + \mu\mat{I} \succ 0$ even if the unregularized $Q_{k,\uu}$
      is singular or ill-conditioned.
\item \textbf{Step-size control:} Large $\mu$ gives conservative (small) steps; $\mu \to 0$ recovers the Newton step.
\end{enumerate}

In the LM setting, the step direction is:
\begin{equation}
\Delta\uvec = -(\nabla^2 J + \mu\mat{I})^{-1}\nabla J,
\label{eq:ch5:lm-step}
\end{equation}
a blend between the Newton direction (gradient of the Hessian) and the gradient direction (step size $\sim 1/\mu$).

\subsection{Adaptive Regularization Schedule}

A practical strategy is to adjust $\mu$ based on the quality of the forward rollout:
\begin{enumerate}
\item Start with $\mu_0$ (e.g., $\mu_0 = 0.1$).
\item After a successful step (cost decreased): $\mu \gets \mu / 10$ (trust the quadratic model more).
\item After a failed step (cost increased despite line search): $\mu \gets 10\mu$ (reduce trust in the model).
\item If $\mu$ exceeds a threshold (e.g., $\mu > 10^6$), the quadratic model is deemed unreliable; restart
      with a simpler approach (e.g., gradient descent) or declare failure.
\end{enumerate}

\subsection{Expected vs Actual Cost Reduction}

The backward pass predicts the cost reduction based on the quadratic model:
\begin{equation}
\Delta J_{\text{expected}} = -\frac{1}{2}\sum_{k=0}^{T-1} \du_k^{(\text{ff})\T}Q_{k,\uu}\du_k^{(\text{ff})}.
\label{eq:ch5:expected-reduction}
\end{equation}

After rolling out the trajectory with the actual nonlinear dynamics, the actual cost reduction is:
\begin{equation}
\Delta J_{\text{actual}} = J_{\text{new}} - J_{\text{old}}.
\label{eq:ch5:actual-reduction}
\end{equation}

The ratio $\rho = \Delta J_{\text{actual}} / \Delta J_{\text{expected}}$ indicates how well the quadratic model
predicts the nonlinear outcome:
\begin{itemize}
\item $\rho \approx 1$: Quadratic model is accurate; decrease $\mu$.
\item $\rho < 0.1$: Model is poor; increase $\mu$ or re-linearize.
\item $\rho < 0$: Cost increased (bad step); always increase $\mu$ or use line search.
\end{itemize}

\section{Cost Weight Selection}

Choosing the matrices $\mat{Q}$ and $\mat{R}$ is as much art as science. Here are practical guidelines.

\subsection{Physical Normalization (Bryson's Rule)}

\textbf{Bryson's Rule} suggests normalizing each state and control by its maximum acceptable magnitude:
\begin{align}
\tilde\x_k &= \frac{\x_k}{x_{\max}}, \quad \tilde\uvec_k = \frac{\uvec_k}{u_{\max}}, \label{eq:ch5:bryson-norm}\\
Q &= \frac{1}{x_{\max}^2}, \quad R = \frac{1}{u_{\max}^2}. \label{eq:ch5:bryson-weights}
\end{align}

For example, if position should not exceed $1$ meter and torque should not exceed $10$ N$\cdot$m, then:
$Q = 1$, $R = 0.01$. This scaling ensures that the cost function naturally penalizes violations of physical
constraints.

\subsection{Frequency-Domain Interpretation}

For a linearized system with dynamics $\dot\x = \mat{A}\x + \mat{B}\uvec$, the LQR cost $\int_0^\infty (\x^\T\mat{Q}\x + \uvec^\T\mat{R}\uvec) dt$
has a frequency-domain interpretation via Bode plots. The closed-loop bandwidth is roughly proportional to $\sqrt{Q/R}$.
If $Q$ is large relative to $R$, the controller is aggressive and has high bandwidth. Conversely, large $R$ gives
a slow, conservative controller.

For the inverted pendulum, the natural frequency of small oscillations around the upright is $\omega_n \approx \sqrt{g/\ell} \approx 3.16$ rad/s.
To balance with a bandwidth slightly faster than this (e.g., $\omega_c = 5$ rad/s), one might choose $Q/R$ to achieve
a closed-loop pole at $-5$ rad/s.

\subsection{Tuning in Practice}

Start with Bryson's rule. If the resulting control is too aggressive or too sluggish:
\begin{itemize}
\item \textbf{Too slow:} Increase $Q$ (penalize errors more) or decrease $R$ (make control cheaper).
\item \textbf{Too aggressive:} Decrease $Q$ or increase $R$ (penalize control effort).
\item \textbf{Oscillatory:} Increase $R$ to damp the response.
\item \textbf{Asymmetric performance:} Use diagonal $\mat{Q}$ with different weights for different states.
\end{itemize}

Systematic tuning can be done by grid search or by tools like the Ziegler--Nichols method (adapted to LQR).
Modern approaches use iterative adjustment based on simulation or hardware experiments.

\section{iLQR vs Full DDP: Computational Trade-offs}
\label{sec:ch5:ilqr-vs-ddp}

\subsection{Computational Complexity}

\textbf{Full DDP} requires:
\begin{itemize}
\item Forward pass: $O(n)$ (integrate nonlinear dynamics).
\item Linearization: $O(n^2 m)$ (Jacobians are $n \times m$ matrices, one per step).
\item Quadratization: $O(n^2 + m^2)$ (Hessians are $n \times n$ and $m \times m$ matrices).
\item Backward pass: Each Riccati step requires inverting $Q_{k,\uu}$ ($O(m^3)$) and updating $\mat{S}_k$ ($O(n^3)$).
  Total: $O(T(n^3 + m^3))$.
\end{itemize}

\textbf{iLQR} avoids the cross-coupling terms:
\begin{itemize}
\item Backward pass: Invert $Q_{k,\uu}$ ($O(m^3)$) but Riccati update is simpler. Total: $O(T \cdot m^3)$ (typically faster
  because $m \leq n$).
\end{itemize}

For typical problems ($n = 10$ states, $m = 2$ controls, $T = 100$ steps):
\begin{itemize}
\item Full DDP: $\approx 100 \times 1000 = 10^5$ floating-point operations per iteration.
\item iLQR: $\approx 100 \times 8 = 800$ floating-point operations per iteration.
\end{itemize}

iLQR is roughly $100 \times$ faster per iteration. However, if iLQR requires $5 \times$ more iterations to converge,
the total time can be comparable. Near the optimum, full DDP's quadratic convergence often wins overall.

\subsection{When Second-Order Terms Matter}

Second-order dynamics terms ($\mat{B}_k^\T\mat{S}_{k+1}\mat{A}_k$) become important when:
\begin{enumerate}
\item The system is underactuated (many states, few controls), so state errors couple strongly through the dynamics.
\item The optimal trajectory involves rapid changes in direction or high accelerations.
\item The cost has explicit cross-coupling ($\ell_{\xu} \neq 0$).
\item The time horizon is long, and cost propagation through the dynamics is significant.
\end{enumerate}

For smooth trajectories with modest time horizons and good actuator bandwidth, iLQR usually suffices.

\section{Continuous-Time Formulation}

\subsection{Hamilton-Jacobi-Bellman Equation}

In continuous time, Bellman's principle becomes the \emph{Hamilton-Jacobi-Bellman (HJB)} equation. Let
$V(t,\x)$ be the value function (optimal cost-to-go) at time $t$ with state $\x(t)$. The HJB equation is:
\begin{equation}
-\frac{\partial V}{\partial t} = \min_{\uvec} \left[\ell(t,\x,\uvec) + \frac{\partial V}{\partial \x}^\T f(t,\x,\uvec)\right].
\label{eq:ch5:hjb}
\end{equation}

For the finite-horizon problem with terminal cost $\ell_T(\x(T))$, the boundary condition is:
\begin{equation}
V(T,\x(T)) = \ell_T(\x(T)).
\label{eq:ch5:hjb-boundary}
\end{equation}

This partial differential equation is difficult to solve in general. However, when the cost and dynamics are
quadratic, the value function is also quadratic, $V(t,\x) = \frac{1}{2}\x^\T\mat{S}(t)\x$, and the HJB equation
reduces to the differential Riccati equation.

\subsection{Differential Riccati Equation (DRE)}

For the continuous-time linear-quadratic regulator with dynamics $\dot\x = \mat{A}(t)\x + \mat{B}(t)\uvec$
and cost
\begin{equation}
J = \frac{1}{2}\x(T)^\T\mat{Q}_T\x(T) + \int_0^T \frac{1}{2}(\x^\T\mat{Q}(t)\x + \uvec^\T\mat{R}(t)\uvec) dt,
\label{eq:ch5:ct-cost}
\end{equation}
the optimal feedback gain is $\uvec^*(t) = -\mat{K}(t)\x(t)$ with $\mat{K}(t) = \mat{R}(t)^{-1}\mat{B}(t)^\T\mat{S}(t)$,
and the cost-to-go matrix $\mat{S}(t)$ satisfies the \textbf{differential Riccati equation}:
\begin{equation}
-\dot{\mat{S}} = \mat{Q} + \mat{A}^\T\mat{S} + \mat{S}\mat{A} - \mat{S}\mat{B}\mat{R}^{-1}\mat{B}^\T\mat{S}.
\label{eq:ch5:DRE}
\end{equation}

Boundary condition (at the terminal time $t = T$): $\mat{S}(T) = \mat{Q}_T$.

The negative sign and backward time direction reflect the fact that $\mat{S}(t)$ evolves \emph{backward} in time from
the terminal cost.

\subsection{Algebraic Riccati Equation (ARE)}

For time-invariant systems and an infinite time horizon, $\mat{S}(t)$ converges to a steady state $\mat{S}_\infty$
satisfying the \textbf{algebraic Riccati equation}:
\begin{equation}
0 = \mat{Q} + \mat{A}^\T\mat{S} + \mat{S}\mat{A} - \mat{S}\mat{B}\mat{R}^{-1}\mat{B}^\T\mat{S}.
\label{eq:ch5:ARE}
\end{equation}

The stabilizability and detectability conditions ensure that a unique positive-definite solution exists, and the
resulting feedback $\mat{K} = \mat{R}^{-1}\mat{B}^\T\mat{S}$ stabilizes the system (renders $\mat{A} - \mat{B}\mat{K}$
Hurwitz).

\subsection{Pontryagin's Minimum Principle in Continuous Time}

The continuous-time analog of the discrete KKT conditions is given by Pontryagin's minimum principle. Define the
Hamiltonian:
\begin{equation}
H(t,\x,\uvec,\vec{\lambda}) = \ell(t,\x,\uvec) + \vec{\lambda}^\T f(t,\x,\uvec).
\label{eq:ch5:hamiltonian-ct}
\end{equation}

Pontryagin's principle states that the optimal control minimizes the Hamiltonian:
\begin{equation}
\uvec^*(t) = \argmin_{\uvec} H(t,\x^*(t),\uvec,\vec{\lambda}(t)),
\label{eq:ch5:minimum-principle}
\end{equation}
and the costates evolve backward via:
\begin{equation}
-\dot{\vec{\lambda}} = \frac{\partial H}{\partial \x} = \frac{\partial \ell}{\partial \x} + \left(\frac{\partial f}{\partial \x}\right)^\T\vec{\lambda}.
\label{eq:ch5:costate-ode}
\end{equation}

For the quadratic case, $\vec{\lambda}(t) = \mat{S}(t)\x(t)$ is proportional to the Riccati solution, and
Pontryagin's conditions recover the standard LQR results.

\begin{intuitionbox}
The HJB equation, DRE, ARE, Pontryagin's principle, and Newton's method applied to the KKT conditions
are all equivalent formulations of the same optimization problem---each offering a different perspective.
The DRE is the most convenient for numerical integration. The ARE is the gateway to classical control
design (pole placement, robust control). Pontryagin's conditions are essential for handling constraints.
Newton's method viewpoint explains why DDP converges so rapidly. Together, they form a unified framework
for understanding optimal control.
\end{intuitionbox}

\section{Practical Considerations}

\subsection{Regularization Revisited}

When $Q_{k,\uu}$ is ill-conditioned or indefinite (which can happen far from optimality), add a regularization term:
\begin{equation}
\mat{K}_k = \bigl(Q_{k,\uu} + \mu\mat{I}\bigr)^{-1}Q_{k,\xu}^\T.
\label{eq:ch5:regularized-gain}
\end{equation}

The parameter $\mu > 0$ acts as a Levenberg--Marquardt damping: large $\mu$ produces conservative (small) steps;
$\mu \to 0$ recovers the full Newton step.

\subsection{Line Search}

After computing the feedback law, the forward rollout uses a step-size parameter $\alpha \in (0,1]$:
\begin{equation}
\uvec_k = \bar\uvec_k + \alpha\,\du_k^{(\text{ff})} + \mat{K}_k(\x_k - \bar\x_k),
\label{eq:ch5:line-search}
\end{equation}
where $\du_k^{(\text{ff})}$ is the feedforward correction. Backtracking line search ensures sufficient cost decrease
(Armijo condition).

\subsection{Convergence Properties}

Near a local optimum where the second-order sufficient conditions hold:
\begin{itemize}
\item DDP with full step ($\alpha = 1$) converges \textbf{quadratically}.
\item iLQR (which omits second-order dynamics terms) converges \textbf{superlinearly}.
\item Both re-linearize at each iteration, ensuring the tangent-space model remains fresh.
\end{itemize}

\section{Model Predictive Control: Online Trajectory Optimization}

DDP and iLQR as presented above are \emph{offline} algorithms: they compute an optimal
trajectory and feedback policy before execution. In practice, disturbances, model mismatch,
and changing objectives require \emph{online} re-optimization. Model Predictive Control
(MPC) provides this capability by embedding the DDP/iLQR computation within a real-time
feedback loop.

\subsection{The Receding-Horizon Principle}

At each control cycle (time $t_k$), the MPC controller:
\begin{enumerate}
\item Measures the current state $\x_k$.
\item Solves a finite-horizon optimal control problem of length $N$:
\begin{equation}
\min_{\uvec_k, \ldots, \uvec_{k+N-1}} \sum_{j=0}^{N-1} \ell(\x_{k+j}, \uvec_{k+j})
+ V_f(\x_{k+N}),
\label{eq:ch5:mpc-problem}
\end{equation}
subject to the dynamics $\x_{k+j+1} = f(\x_{k+j}, \uvec_{k+j})$.
\item Applies only the first control $\uvec_k^*$ to the system.
\item Shifts the horizon forward and repeats at $t_{k+1}$.
\end{enumerate}

The key insight is that MPC is precisely DDP/iLQR executed at each time step, with the
previous solution shifted forward as the initial guess (warm-starting).

\subsection{Warm-Starting and Computational Efficiency}

When the system evolves smoothly, the optimal trajectory at time $t_{k+1}$ is close to
the shifted trajectory from time $t_k$. Concretely, if $\{\bar\x_j, \bar\uvec_j\}_{j=k}^{k+N}$
was optimal at $t_k$, the warm-start for $t_{k+1}$ is:
\begin{equation}
\bar\x_j^{(0)} = \bar\x_{j+1}^{\text{prev}}, \quad
\bar\uvec_j^{(0)} = \bar\uvec_{j+1}^{\text{prev}}, \quad j = k+1, \ldots, k+N-1,
\label{eq:ch5:warm-start}
\end{equation}
with the final element filled by a reasonable extrapolation.

This warm-starting exploits the continuity of the value function and typically reduces the
number of DDP iterations from many (cold start) to one or two (warm start).

\subsection{Terminal Cost and Stability}

The terminal cost $V_f(\x_{k+N})$ in~\eqref{eq:ch5:mpc-problem} plays a critical role
in guaranteeing stability. The standard approach is to choose $V_f$ as the \emph{infinite-horizon
LQR cost} at the target:
\begin{equation}
V_f(\x) = \x^T \mat{S}_\infty \x,
\label{eq:ch5:terminal-cost}
\end{equation}
where $\mat{S}_\infty$ is the solution of the algebraic Riccati equation at the target
linearization. Under this choice:

\begin{theorem}[MPC Stability]
\label{thm:ch5:mpc-stability}
If $V_f$ is a local Control Lyapunov\index{Lyapunov Stability} Function (CLF) for the system at the target
equilibrium, and the terminal constraint $\x_{k+N} \in \mathcal{X}_f$ (a sublevel set of
$V_f$) is imposed, then the MPC law is asymptotically stabilizing.
\end{theorem}

The connection to contraction theory (Chapter~\ref{ch:contraction}) is direct: the LQR
terminal cost provides a contraction metric (Chapter~\ref{ch:duality}) that guarantees
exponential convergence within the terminal set. MPC stability is therefore a consequence
of the stability--optimality duality.

\subsection{The Real-Time Iteration (RTI) Scheme}

When computational resources are limited (e.g., embedded systems), the RTI approach
executes exactly \emph{one} Newton step (one backward Riccati pass + one forward rollout)
per control cycle, rather than iterating DDP to convergence. The RTI scheme:

\begin{enumerate}
\item Linearizes the dynamics and quadraticizes the cost around the current shifted trajectory.
\item Performs one backward Riccati sweep to compute $\mat{K}_k$ and $\vec{k}_k$.
\item Applies $\uvec_k = \bar\uvec_k + \vec{k}_k + \mat{K}_k(\x_k - \bar\x_k)$.
\item Shifts the trajectory and prepares for the next cycle.
\end{enumerate}

This is computationally equivalent to applying a single-iteration DDP at each time step.
The convergence analysis shows that if the system changes slowly relative to the control
rate, the RTI scheme tracks the optimal solution with bounded error.

\begin{designbox}
MPC unifies the offline trajectory optimization of this chapter with the online feedback
requirements of real systems. The computational loop---linearize, solve Riccati, apply
control, re-linearize---is exactly the tangent-space methodology of this book, executed
in real time. Modern implementations on embedded hardware achieve cycle times of 1--10~ms
for systems with $n \leq 50$ states and $N \leq 100$ horizon steps, making MPC practical
for robotics, automotive, and aerospace applications.
\end{designbox}


\section{Chapter Summary}

\begin{enumerate}
\item \textbf{Historical evolution:} From Bellman's dynamic programming (1957) to Pontryagin's maximum principle (1962)
      to DDP (1970), optimal control theory evolved from multiple independent viewpoints that converge on a unified
      framework.

\item \textbf{Local quadratic subproblem:} Nonlinear costs and dynamics are expanded to second order around a nominal
      trajectory, yielding the Q-function with Hessian terms $Q_{k,\xx}$, $Q_{k,\uu}$, $Q_{k,\xu}$.

\item \textbf{Riccati recursion:} Solves the local quadratic subproblem exactly, computing both the optimal feedback
      gain $\mat{K}_k$ and the cost-to-go matrix $\mat{S}_k$. iLQR simplifies by dropping cross-coupling; full DDP
      retains it for higher accuracy.

\item \textbf{DDP algorithm:} Iterates the Riccati recursion by re-linearizing and re-quadratizing at each trajectory
      update. This is equivalent to Newton's method applied to the trajectory optimization KKT conditions.

\item \textbf{Convergence:} DDP exhibits local quadratic convergence near a local optimum, inheriting Newton's superb
      final convergence rate. Line search and regularization ensure global descent away from the optimum.

\item \textbf{Worked example:} Inverted pendulum swing-up demonstrates cost reduction over iterations and the structure
      of feedback gains for an underactuated system.

\item \textbf{Constraints:} Control bounds via clamping, state constraints via penalties or barriers, and trust regions
      are practical approaches to enforce constraints.

\item \textbf{Cost weighting:} Bryson's rule normalizes by maximum acceptable deviations. Frequency-domain interpretation
      links weights to closed-loop bandwidth. Iterative tuning refines performance.

\item \textbf{Computational trade-offs:} Full DDP is $O(T n^3)$ but exhibits quadratic convergence. iLQR is $O(T m^3)$
      and often faster overall, though with superlinear rather than quadratic convergence. Choice depends on problem
      structure and required accuracy.

\item \textbf{Continuous time:} HJB equation, DRE, and ARE are the continuous-time analogs. Pontryagin's principle
      extends optimality conditions to constrained problems. The ARE is the foundation for classical control design.

\item \textbf{Duality:} The cost-to-go matrix $\mat{S}_k$ has dual interpretations as cost curvature and as a candidate
      contraction metric---a connection formalized in the next chapter on contraction analysis.
\end{enumerate}

\begin{axiombox}
\textbf{Principle of Local Optimality:} The trajectory optimization problem, though globally nonlinear and nonconvex,
can be solved iteratively by working in local tangent spaces. At each step, we solve an exact, finite-dimensional,
convex subproblem (the LQR problem). By re-centering the tangent space at each iteration, we avoid committing to a
global approximation. This principle underlies not only DDP but also many modern algorithms in robotics and machine
learning (e.g., SAC, MPPI, trajectory optimization in MPC).
\end{axiombox}

\section*{Exercises}

\begin{enumerate}
\item \textbf{LQR for a double integrator.}
      Solve the discrete-time LQR problem for $\x_{k+1} = \begin{bmatrix} 1 & h \\ 0 & 1 \end{bmatrix}
      \x_k + \begin{bmatrix} 0 \\ h \end{bmatrix} u_k$ with $\mat{Q} = \mat{I}_2$, $R = 1$,
      $h = 0.01$, and horizon $N = 100$. Plot the trajectory and gain $\mat{K}_k$.

\item \textbf{Q-function structure.}
      For the inverted pendulum from the worked example, compute $Q_{k,\xx}$, $Q_{k,\uu}$,
      and $Q_{k,\xu}$ at the first step. Verify that $Q_{k,\uu} \succ 0$ and interpret the
      eigenvalues of $Q_{k,\xx}$.

\item \textbf{iLQR vs full DDP.}
      Implement both iLQR and full DDP for the cart-pole system. Run each for 50 iterations
      starting from the same initial guess. Compare convergence rates and final costs.
      When do they differ significantly?

\item \textbf{Regularization effects.}
      For the inverted pendulum, run DDP with regularization parameters
      $\mu \in \{0, 0.01, 0.1, 1, 10\}$. Plot convergence (cost vs. iteration) for each.
      Explain the trade-off between step size and convergence rate.

\item \textbf{MPC horizon sensitivity.}
      Implement MPC for the inverted pendulum with horizons $N \in \{10, 20, 50, 100\}$.
      Use the infinite-horizon LQR cost as terminal cost. Compare tracking performance
      under a disturbance at $t = 1$ s. How short can $N$ be before performance degrades?

\item \textbf{Cost weight selection.}
      For the pendulum swing-up, experiment with different $\mat{Q}/\mat{R}$ ratios.
      Plot the resulting trajectories and control efforts. Explain why very large $\mat{Q}$
      can cause DDP to diverge (relate to the condition number of $Q_{k,\uu}$).

\item \textbf{Warm-starting effectiveness.}
      Implement MPC with and without warm-starting. Count the number of DDP iterations
      needed per time step in each case. Plot the computational savings as a function
      of the control rate.

\item \textbf{Continuous-time HJB.}
      For the scalar system $\dot{x} = -x + u$ with cost $\int_0^\infty (x^2 + u^2)\,dt$,
      solve the continuous-time ARE analytically. Verify that $V(x) = Sx^2$ satisfies the
      HJB equation $0 = x^2 + Sx^2/S - Sx\cdot x + Sx(-x + u^*)$.
\end{enumerate}
